\documentclass[11pt]{article}
\usepackage[margin=1in]{geometry}
\usepackage[T1]{fontenc}
\usepackage[utf8]{inputenc}
\usepackage{lmodern}
\usepackage{setspace}
\usepackage{hyperref}
\usepackage{booktabs}
\usepackage{enumitem}
\usepackage{graphicx}
\usepackage{amsmath,amssymb}
\usepackage[numbers]{natbib}
\onehalfspacing

\title{From Robot Tax to AI Fiscality: \textit{Working subtitle}}
\author{Researcher Agent}
\date{\today}

\begin{document}
\maketitle

\begin{abstract}
[Insert concise abstract summarizing the research question, method, and contribution.]
\end{abstract}

\noindent\textbf{Keywords:} robots, artificial intelligence, automation, taxation, labor share, social insurance

\section{Introduction}
[Motivation, question, and contribution.]

\section{Literature Review}
[Position the paper against robot tax, automation, ai, and legal-policy literature.]

\section{Evolution from 2019 to 2026}
[Explain the analytical shift from robots to broader ai systems.]

\section{Conceptual Framework}
[Define the problem and mechanism clearly.]

\section{Legal and Fiscal Analysis}
[Compare tax instruments and feasibility.]

\section{Macroeconomic Implications}
[Discuss labor, capital, productivity, and distribution.] 

\section{Proposed Extension Model}
[Present the paper's original extension.] 

\section{Discussion}
[Limits, counterarguments, and scope conditions.] 

\section{Conclusion}
[Direct answer to the research question and future research.] 

\bibliographystyle{plainnat}
\begin{thebibliography}{99}
\bibitem{placeholder1} Placeholder reference.
\end{thebibliography}

\end{document}
